\section{Введение}
\label{sec:intro}

Специализированные ускорители нейросетей (GPU, TPU, NPU) оптимизированы под
ключевую операцию нейросетевых вычислений~--- обобщённое матричное умножение
(General Matrix Multiply, GEMM). Разработка программного обеспечения для таких
устройств ведётся с помощью \textbf{функциональных симуляторов}~--- программных
моделей устройства, исполняемых на стандартном оборудовании и воспроизводящих
его поведение на уровне отдельных операций. Ключевое требование к симулятору~---
соответствие его числовых результатов результатам целевого устройства.

Поскольку GEMM доминирует в вычислительном графе, именно её производительность
определяет применимость симулятора. Последовательная CPU-реализация GEMM
неприемлемо медленна, поэтому вычисление переносят на GPU с помощью библиотеки
cuBLAS~\cite{cuBLAS2024} или специализированных ядер~\cite{CUDA2024}. Такая
замена неизбежно изменяет числовой результат: сложение в формате
\texttt{FP32} неассоциативно~\cite{Higham2002,Goldberg1991}, и каждая схема
накопления частичных произведений (последовательная на CPU, блочная на GPU,
систолическая в ускорителе) порождает свою последовательность ошибок округления.

Влияние расхождений не ограничивается одной операцией. Аппаратная ошибка слоя
переносится как возмущение на вход следующего, усиливаясь весовыми матрицами с
большой спектральной нормой и подавляясь нормализационными слоями. При инференсе
расхождения накапливаются вдоль прямого прохода, при обучении~--- искажают
градиенты обратного прохода. При этом характер накопления принципиально зависит
от природы ошибки: случайные расхождения частично компенсируются при
суммировании по слоям, тогда как систематически смещённые складываются
конструктивно~\cite{Higham2002}. Возникает структурное противоречие
разработки симулятора: выбор реализации GEMM влияет и на скорость, и на
достоверность результатов.

Настоящая работа исследует влияние выбора реализации GEMM в функциональном
симуляторе на вычислительную точность при инференсе и обучении. Рассматриваются
три реализации~--- последовательная CPU-реализация (точка отсчёта), cuBLAS и
специализированное CUDA-ядро. Исследование выполнено на функциональном симуляторе
промышленного специализированного ускорителя; конкретное устройство и симулятор
не раскрываются в силу соглашения о неразглашении, однако предложенные модель и
измерительный конвейер не зависят от целевой платформы и применимы к любому
функциональному симулятору с заменяемой реализацией GEMM. \textbf{Вклад работы:}

\begin{enumerate}\setlength\itemsep{2pt}
    \item аналитическая модель накопления расхождений по слоям через матрицы
        Якоби; установлена роль слоёв со стоками ($\|J\|_2<1$~--- нормализация,
        насыщающиеся нелинейности, усреднения), раздел~\ref{sec:model};
    \item измерительный конвейер с раздельной оценкой вклада компилятора и
        реализации GEMM и диагностикой типа ошибки за один прямой проход,
        раздел~\ref{sec:method};
    \item эмпирический вывод: доминирует компилятор XLA, поэтому относительно
        CPU-эталона перенос GEMM на GPU не ухудшает точность ($R\le1.7$ для
        случайной ошибки в прямом проходе), раздел~\ref{sec:results};
    \item тип ошибки (смещённая или случайная) служит более информативным
        предиктором накопления, чем её поэлементная норма.
\end{enumerate}

Изложение построено следующим образом. Раздел~\ref{sec:related} обозревает
связанные работы; раздел~\ref{sec:model} вводит модель накопления;
раздел~\ref{sec:method} описывает измерительный конвейер;
раздел~\ref{sec:results} приводит результаты; разделы~\ref{sec:discussion}
и~\ref{sec:conclusion} содержат обсуждение и заключение.
